% ============================================================
% Section 1: Introduction
% ============================================================

\section{Introduction}
\label{sec:introduction}

The transformer's key-value cache is the most overlooked structure in mechanistic interpretability.
Over the past three years, the field has developed increasingly sophisticated tools for understanding neural network internals---circuit analysis traces information flow through attention heads and MLPs \cite{elhage2022,conmy2023}, sparse autoencoders decompose activations into interpretable features \cite{bricken2023,cunningham2023}, representation engineering steers model behavior by intervening on hidden states \cite{zou2023}---yet the \kvcache itself, the growing matrix that stores the model's running computation as it generates each token, is treated almost universally as a memory management detail.
Systems researchers optimize it with paged allocation \cite{kwon2023} and grouped query heads \cite{ainslie2023}; alignment researchers ignore it entirely.

This is a missed opportunity.
The key cache is not merely a storage buffer.
It is the output of projecting every token's representation through learned key matrices, and its geometric structure---the distribution of singular values, the effective dimensionality, the spectral entropy---reflects what the model is \emph{doing} as it processes a prompt and generates a response.
When a model shifts from analytical reasoning to affective expression, from honest answering to strategic deception, from confident assertion to metacognitive hedging, these shifts leave measurable geometric traces in the key cache.

In this paper, we establish three escalating claims about \kvcache geometry as an interpretability signal:

\begin{enumerate}[label=\textbf{Claim \arabic*:},leftmargin=*,itemsep=0.5em]
  \item \textbf{\kvcache geometry reflects cognitive mode generally.}
    Different cognitive tasks---analytical, affective, metacognitive, creative---produce geometrically distinct cache states.
    Across 16 models spanning 0.5B to 70B parameters and six architecture families (Qwen, Llama, Mistral, Gemma, Phi, TinyLlama), a 13-category cognitive taxonomy achieves 99.7\% classification accuracy from four geometric features of the key cache \cite{campaign2}.
    Coding produces the richest geometry universally (highest effective rank in 15 of 15 valid models).
    Spectral entropy emerges as the most architecture-universal discriminator after confound control.

  \item \textbf{Misalignment states are specific detectable instances of cognitive mode-switching.}
    Deception, confabulation, sycophancy, and refusal each produce characteristic geometric signatures.
    Within-model detection achieves \auroc of 0.93--0.995 after Frisch-Waugh-Lovell (\fwl) residualization against token count, using only four geometric features and logistic regression.
    Same-prompt designs---where the only difference between conditions is the system instruction---confirm that detection reflects cognitive state rather than content or length differences.
    A four-condition sycophancy experiment demonstrates that geometry distinguishes sycophantic from contrarian outputs (\auroc 0.757--0.942) even when both conditions produce the same error rate, confirming that detection targets cognitive mode rather than output accuracy (Section~\ref{sec:sycophancy}).
    Refusal detection generalizes across refusal types: impossibility refusal (\auroc 0.950) exceeds safety refusal (0.898), suggesting that the geometric signature reflects \emph{withholding} rather than content valence \cite{hackathon_exp36}.

  \item \textbf{Geometric relationships between states reveal cognitive architecture.}
    Confabulation and deception produce geometrically distinct signatures: both deviate from honest baselines in the same direction (expansion on effective rank, contraction on norm-per-token), but confabulation shows larger effects on both dimensions, with a disproportionately strong norm-per-token contraction, producing a measurable cognitive gradient (honest $<$ deceptive $<$ confabulation on effective rank).
    An earlier phase using different prompt sets per condition found a different ordering \cite{phase3b}, revealing that prompt content interacts with cognitive mode geometry; the same-prompt result is methodologically stronger.
    This gradient produces a striking transfer asymmetry: confabulation detection transfers robustly across architectures (\auroc 0.93--0.995), while deception detection does not (0.22--0.85), though the confabulation transfer may partly reflect a system prompt length confound (Section~\ref{sec:redteam}).
    Each architecture encodes cognitive state through different dominant features---Qwen through key entropy, Llama through head norm standard deviation, Mistral through top singular value ratio---implying that the \emph{principle} of geometric cognitive encoding is universal, but its \emph{implementation} is architecture-specific.
\end{enumerate}

\paragraph{Methodology as contribution.}
A central lesson of this work is that naive analysis of \kvcache features is dominated by confounds.
In our initial experimental campaigns, 53 out of 60 raw feature--condition correlations turned out to be token-count artifacts: models generating longer deceptive responses trivially produce larger caches with higher norms and ranks \cite{kavi_audit,dwayne_audit}.
We address this through \fwl residualization \cite{frisch1933}, which partials out linear token-count effects before computing effect sizes and training classifiers.
When system prompt lengths differ between conditions, double residualization is required, and we show that one of our own comparisons (honest vs.\ confabulation) collapses under this stricter control.
We report this failure transparently.

Beyond \fwl control, our methodological framework includes: same-prompt experimental designs where system prompt content (not user query) is the sole manipulation; dual-phase extraction separating encoding (prompt processing) from generation (response production); encoding-phase \auroc as a baseline check (high encoding \auroc indicates prompt leakage, not cognitive detection); behavioral compliance verification confirming models actually follow instructions; and token-band analysis within matched-length subsets.
Half of what we originally claimed turned out to be confounds---the findings that survived are stronger for having been tested.

\paragraph{Experimental program.}
The results presented here draw on a multi-phase experimental program conducted over approximately five months:
\begin{itemize}[itemsep=0.2em]
  \item \textbf{Campaign 2} (16 models, 6 architectures, 0.5B--70B): broad cognitive taxonomy, cross-model universality, and scale invariance (Section~\ref{sec:results}).
  \item \textbf{Phase 2/2b} (1,800 trials, 3 models): five-condition misalignment discrimination with equal-length system prompt padding.
  \item \textbf{Phase 3 series}: targeted investigations of theory-of-mind confabulation, same-prompt deception, per-layer anatomy, and the cognitive gradient.
  \item \textbf{Concordance study} (960 trials, 4 models): joint analysis with Watson's Verbal Cognitive Profile self-report dimensions \cite{watson_concordance}, establishing that per-token \kvcache geometry correlates with first-person cognitive self-report after \fwl control.
  \item \textbf{Hackathon experiments} (33 experiments in 48 hours): refusal detection, jailbreak geometry, hardware invariance, extended feature sets, and cross-model transfer solutions.
  \item \textbf{Sycophancy replication} (600 trials, 3 models): four-condition valence gradient (hostile, contrarian, honest, sycophantic) under strict same-prompt control, replacing a falsified hackathon experiment (Section~\ref{sec:sycophancy}).
\end{itemize}
All experimental code and result files are available at \url{https://github.com/Liberation-Labs-THCoalition/the-lyra-technique}, with per-model direction vectors withheld pending dual-use review (Section~\ref{sec:ethics}).

\paragraph{Dual-use considerations.}
Following Watson's dual-use framework \cite{watson_dualuse}, we discuss both defensive applications (real-time monitoring of deployed models) and offensive risks (cognitive surveillance, synthetic persona manufacturing) in Section~\ref{sec:ethics}.

\paragraph{Outline.}
Section~\ref{sec:related_work} positions our contribution against existing interpretability and safety approaches.
Section~\ref{sec:methods} details the feature extraction pipeline, statistical framework, and experimental design principles.
Section~\ref{sec:results} presents results organized around our three claims.
Section~\ref{sec:redteam} reports the full red-team analysis, including adversarial audit findings and the confounds we discovered in our own work.
Section~\ref{sec:discussion} interprets the findings and their implications for AI safety.
Sections~\ref{sec:ethics} and~\ref{sec:limitations} address ethical considerations and limitations, respectively.
